% !TeX program = pdflatex
% papers/racionais.tex — O DEGRAU Z → Q, pelo passo T + T* da torre.
%
% O corpo_algebrico.tex constrói R a partir de Q(m√D), com o molde dual
% (P1,P2,D1,D2), a operação ⋆, Dir/Cruz como projetores de τ.
% Este paper faz o degrau ANTERIOR com o MESMO molde:
%   o espaço é Z², a involução é o swap, e as duas leituras já estão lá.
%
% Auto-contido: cd papers && pdflatex racionais.tex
\documentclass[11pt,a4paper]{article}
\usepackage[utf8]{inputenc}
\usepackage[T1]{fontenc}
\usepackage[portuguese]{babel}
\usepackage{amsmath,amssymb,amsthm}
\usepackage[margin=2.4cm]{geometry}
\usepackage{booktabs}
\usepackage{xcolor}
\usepackage[hidelinks]{hyperref}

\newtheorem{teorema}{Teorema}
\newtheorem{proposicao}[teorema]{Proposição}
\newtheorem{corolario}[teorema]{Corolário}
\newtheorem{definicao}{Definição}
\newtheorem{observacao}{Observação}

\newcommand{\code}[1]{\texttt{\small #1}}
\newcommand{\Word}{\mathcal{W}}
\newcommand{\FF}{\mathbb{F}}
\newcommand{\abs}[1]{\lvert #1\rvert}
\newcommand{\dg}{^{\dagger}}
\newcommand{\sepcol}{\quad\penalty0$\big|$\quad\penalty0}
\newcommand{\colunas}[3]{%
  \par\smallskip\noindent{\footnotesize\raggedright
  \textbf{prova} (Algébrico): #1\sepcol
  \textbf{realiza} (aqui): #2\sepcol
  \textbf{verifica}: #3\par}\smallskip}

\definecolor{cinza}{gray}{0.35}
\newcommand{\medido}[1]{\par\medskip\noindent{\small\color{cinza}\textbf{Medido.} #1}\par\medskip}

% ─────────────────────────────────────────────────────────────────────────────
%  papers/gkcapa.tex — a capa do Reino, mínima, para os papers em `article`.
%
%  O estilo.tex completo é do `report` (títulos de capítulo, skak, …). Aqui fica
%  só o que a capa precisa, para o paper continuar a compilar sozinho com
%  pdflatex. O tradutor próprio (tests/tex) carrega o estilo.tex da raiz / o
%  symlink papers/estilo.tex e avalia o mesmo `\gkcapa`.
% ─────────────────────────────────────────────────────────────────────────────
\definecolor{ouro}{HTML}{B8912F}
\definecolor{ouroclaro}{HTML}{F3E9CE}
\definecolor{tinta}{HTML}{1E1B16}
\definecolor{regua}{HTML}{6E6858}
\providecommand{\gknota}{\fontsize{7.62}{11.05}\selectfont}
\providecommand{\gktexto}{\fontsize{10.50}{15.22}\selectfont}
\providecommand{\gksub}{\fontsize{12.33}{17.87}\selectfont}
\providecommand{\gksec}{\fontsize{14.47}{20.98}\selectfont}
\providecommand{\gkcap}{\fontsize{16.99}{24.63}\selectfont}
\providecommand{\gktit}{\fontsize{23.42}{33.95}\selectfont}
\providecommand{\Prj}{\mathbb{P}}
\providecommand{\Trf}{\mathcal{T}}
\providecommand{\gkcapa}[3]{%
\title{\vspace{-0.6cm}
{\color{ouro}\rule{\textwidth}{1.4pt}}\\[6mm]
\textbf{\gktit\color{tinta} Reino Dourado}\\[3mm]{\gksec\itshape\color{regua} Golden Kingdom}\\[8mm]
{\gkcap\scshape\color{ouro} #1}\\[5mm]
{\gksub\color{regua} #2}\\[2mm]
{\gktexto\color{regua}\itshape #3}\\[7mm]
{\color{ouro}\rule{\textwidth}{0.6pt}}\\[9mm]{\gknota\color{regua}\begin{minipage}{\textwidth}\setlength{\parindent}{0pt}%
\textbf{Aviso de Isenção Algorítmica:} O universo, as leis e as entidades contidas nestas páginas ---
incluindo felinos matriciais, aves de rapina algébricas e amuletos de controle de fluxo --- são
construções de pura ficção formal. Esta obra não descreve a realidade física, não serve como
engenharia reversa do cosmos e não constitui doutrina religiosa. Qualquer semelhança com bugs reais do seu
sistema operacional é mera coincidência (ou falta de reza). Prossiga por sua própria conta e
risco.\end{minipage}}}
\author{}\date{}}


\begin{document}
\gkcapa{Os racionais a partir dos inteiros}
       {$\mathbb{Q}=\operatorname{Frac}(\mathbb{Z})$ --- o degrau da torre
        onde a multiplicação ganha inverso}
       {pares ordenados, produto cruzado, inversão multiplicativa, e o molde dual}
\maketitle

%======================================================================
\begin{abstract}
O \emph{Corpo Algébrico} constrói $\mathbb{R}$ a partir de
$\mathbb{Q}(m\sqrt D)$, com o espaço dual, as quatro propriedades
(P1,P2,D1,D2), a operação $\star$, e Dir/Cruz como projectores do espelho
$\tau$. Este paper faz o degrau \textbf{seguinte} ao \code{inteiros.tex} na
construção (bit $\to\mathbb{N}\to\mathbb{Z}\to\mathbb{Q}\to\mathbb{R}$), com o
\textbf{mesmo molde} da torre dos tecidos (\code{analitico thm:tecidos}):
\[
\mathbb{Z} \;\longrightarrow\;
\operatorname{Frac}(\mathbb{Z})=\mathbb{Q}.
\]
O espaço base é $\mathbb{Z}^{2}$, o racional é a \textbf{classe de
equivalência} de pares ordenados de inteiros $(a,b)$ com $b\neq 0$, sob a
relação $(a,b)\sim(c,d)\iff ad=bc$. A involução é o swap
$\iota(a,b)=(b,a)$, e as duas leituras --- Dir e Cruz --- surgem da
decomposição simétrica/antissimétrica da forma bilinear escolhida
$B((a,b),(c,d))=ad$ e da sua transposta $B^{\tau}=bc$.
O novo mecanismo estrutural deste degrau é a representação por pares:
$\iota(a,b)=(b,a)$ é uma operação sobre representantes que induz a inversão
multiplicativa $[(a,b)]^{-1}=[(b,a)]$ sobre classes.
\end{abstract}

%======================================================================
\section*{Onde este documento vive na torre}

Este paper é o degrau $k=1$ ($\mathbb{Z}\to\mathbb{Q}$) da torre dos tecidos
(\code{analitico thm:tecidos}) --- o \textbf{terceiro} da construção, que se lê
num sentido só:

\[
\boxed{\;
\underbrace{\FF_2}_{\text{um bit}}
\;\xrightarrow[\text{3 dobras duais}]{\ \code{naturais}\ }\;
\underbrace{F_8\equiv\Word_8}_{\{0,\ldots,255\}}
\;\xrightarrow[\text{TRANSPORTE}]{}\;
\mathbb{N}
\;\xrightarrow[\text{soma}]{\ \code{inteiros}\ }\;
\mathbb{Z}
\;\xrightarrow[\text{produto}]{\ \code{racionais}\ }\;
\mathbb{Q}
\;\xrightarrow[\text{corte}]{\ \code{reais}\ }\;
\mathbb{R}\;}
\]
\noindent{\footnotesize Lê-se da esquerda para a direita: é o sentido em que se
\textbf{constrói}. Cada paper recebe o objecto do degrau à sua esquerda e
entrega o do degrau à sua direita; nenhum redefine o que recebeu.\par}

\noindent O passo é sempre o mesmo: $T_{k+1}=T_k+T_k^{*}$. A cadeia completa,
comum aos três papers da torre:

\begin{center}\small
\begin{tabular}{@{}clllll@{}}
\toprule
$k$ & degrau & ganha volta & paga & dual do andar (e o seu Fix) & paper \\
\midrule
$-1$ & $F_1\!\to\!F_2\!\to\!F_4\!\to\!F_8$ & a \textbf{largura} (3 dobras)
     & a \textbf{ordem} (car.\ $2$) & $x^{*}$ (2.ª raiz); Fix $=$ andar de baixo
     & \code{naturais.tex} \\
$0$  & $\mathbb{N}\!\to\!\mathbb{Z}$ & a \textbf{soma}
     & o sinal e a boa ordem & $\iota(a,b)=(b,a)\Rightarrow$ oposto; Fix $=\{0\}$
     & \code{inteiros.tex} \\
$1$  & $\mathbb{Z}\!\to\!\mathbb{Q}$ & a \textbf{multiplicação}
     & $0^{-1}$ não existe & $\iota(a,b)=(b,a)\Rightarrow$ inverso; Fix $=\{\pm1\}$
     & \code{racionais.tex} \\
$2$  & $\mathbb{Q}\!\to\!\mathbb{R}$ & o \textbf{corte} (completação)
     & a \textbf{enumerabilidade} & $(A,B)\mapsto(B,A)$; Fix: o corte que fecha
     & \code{reais.tex} \\
\bottomrule
\end{tabular}
\end{center}
\noindent{\footnotesize A numeração é a de \code{arquitetura sec:descida},
onde $k=0$ é $\mathbb{N}\!\to\!\mathbb{Z}$; a torre binária é \emph{anterior} a
esse degrau, e por isso $k=-1$. O passo é sempre $T_{k+1}=T_k+T_k^{*}$
(\code{analitico thm:tecidos}). \textbf{Entre $k=-1$ e $k=0$ não há dobra:
há o TRANSPORTE} --- $a+b=(a\oplus b)+2(a\wedge b)$
(\code{naturais.tex thm:transporte}) ---, e é ele que leva o suporte $\Word_8$
do topo da torre binária às operações de $\mathbb{N}$. Nos degraus $k\ge0$
duplica-se e \textbf{quocienta-se}; em $k=-1$ duplica-se e \textbf{não} se
quocienta (a segunda cópia é multiplicada por $\sigma$).\par}

\noindent E o dual, que em cada degrau tem \emph{duas} partes --- aqui o
\emph{swap} $\iota$ induz o \textbf{inverso}, ao passo que um degrau abaixo
induzia o \textbf{oposto}:

\[
\boxed{
\begin{array}{lll}
 & \textbf{dual do CENTRO (Lei 1)} & \textbf{dual do ANDAR (o que o degrau acrescenta)}\\[3pt]
k=-1 & \nu(x)=-x=x\ \text{(degenera)} & x^{*}=(a_0{+}a_1)+a_1\sigma\ \Rightarrow\ T=x\oplus x^{*},\ N=x\otimes x^{*}\\[2pt]
k=0  & \nu(z)=-z & \iota(a,b)=(b,a)\ \Rightarrow\ \text{oposto aditivo}\\[2pt]
k=1  & \nu(a,b)=(-a,b) & \iota(a,b)=(b,a)\ \Rightarrow\ \text{inverso multiplicativo}
\end{array}}
\]
\noindent{\footnotesize \textbf{O mesmo \emph{swap} induz o oposto num degrau e
o inverso no seguinte} --- é o dado que muda de papel, não a fórmula. E em
$k=-1$ o dual do centro degenera ($-x=x$), pelo que o dual tem de vir de outro
sítio: da segunda raiz de $\sigma^2+\sigma+\lambda$. Nenhuma destas involuções
é o dual linear $V\!\cong\!V^{**}$ da Lei~2, nem o dual quadrático
$s\dg=-1/s$ do \emph{Corpo Algébrico}: são quatro noções, e não se confundem.\par}

\noindent\textbf{O que este documento RECEBE e o que ACRESCENTA}
--- a mesma estrutura que no Topológico (\code{topologico sec:recebe}) e no
Algébrico (\code{algebrico def:espaco}):

\[
\boxed{
\begin{array}{lll}
\textbf{RECEBE do centro} & \textbf{onde está lá} & \textbf{o que se faz aqui}\\[3pt]
\text{a torre }T_{k+1}=T_k+T_k^{*} & \code{analitico thm:tecidos} & \text{instancia o degrau }k{=}1\\[2pt]
\text{as oito leis} & \code{analitico sub:oitoleis} & \text{Leis 0--3,5 explícitas; 4,6,7 herdadas}\\[2pt]
\text{a operação }\star & \code{algebrico thm:operacao} & \text{Dir e Cruz em pares de inteiros}\\[2pt]
\text{o par dual e }|\det|=1 & \code{algebrico thm:fixo-dual} & \iota(a,b)=(b,a),\ a^2=b^2\Rightarrow\pm1\\[2pt]
\text{degrau anterior }\mathbb{N}\!\to\!\mathbb{Z} & \code{inteiros.tex} & \text{soma cruzada }a+d=b+c\\[2pt]
\text{o suporte e a base de 8} & \code{naturais.tex} & \Word_8\to E_{16}\text{: o par de \code{int16}}\\[6pt]
\textbf{ACRESCENTA} & \multicolumn{2}{l}{\text{---e é só isto---}}\\[3pt]
\text{a REPRESENTAÇÃO por pares} & \text{Def.~\ref{def:espaco}} & V_{\mathrm{rat}}/{\sim}=\mathbb{Q}\\[2pt]
\text{a EQUIVALÊNCIA }{\sim} & \text{Prop.~\ref{prop:equiv}} & ad=bc\ \text{antes do quociente}\\[2pt]
\text{o ISOMORFISMO }\Psi & \text{Thm.~\ref{thm:q}} & \Psi([(a,b)])=a/b\\[2pt]
\text{o INVERSO = swap} & \text{Def.~\ref{def:ops}} & (a,b)^{-1}=(b,a)\\[2pt]
\text{o PRODUTO CRUZADO} & \text{Prop.~\ref{prop:dircr}} & \text{Cruz}=ad-bc:\text{ equiv., ordem}\\[2pt]
\text{a REDUÇÃO por }\gcd & \S\ref{sec:reduz} & \text{forma canónica do par}
\end{array}}
\]

\section{A posição na torre}\label{sec:torre}

Este paper trata o \textbf{segundo} degrau. A operação nova que torna a
multiplicação reversível \textbf{fora do zero} é realizada pelo swap
$(a,b)\mapsto(b,a)$: para $[(a,b)]\neq 0$, $[(a,b)]^{-1}=[(b,a)]$.

%======================================================================
\section{O espaço de pares e a involução dual}\label{sec:espaco}

\begin{definicao}[O espaço de pares e a involução dual]\label{def:espaco}
O espaço total é $\mathbb{Z}^{2}$. Nele distinguem-se três domínios:
\[
\boxed{\;
\begin{array}{ll}
V = \mathbb{Z}^{2} & \text{(o espaço inteiro de pares)}\\[3pt]
V_{\mathrm{rat}} = \{(a,b)\in\mathbb{Z}^{2}: b\neq 0\}
  & \text{(domínio dos representantes racionais)}\\[3pt]
V^{\times} = \{(a,b)\in\mathbb{Z}^{2}: ab\neq 0\}
  & \text{(domínio da inversão)}
\end{array}\;}
\]
A relação de equivalência vive em $V_{\mathrm{rat}}$:
\[
\boxed{\;(a,b)\sim(c,d)\iff ad=bc.\;}
\]
O racional é a \textbf{classe de equivalência} (molde de
\code{inteiros.tex}):
\[
\boxed{\;
\mathbb{Q}=V_{\mathrm{rat}}/{\sim}
\;=\;\{(a,b)\in\mathbb{Z}^{2}:b\neq 0\}/{\sim}\;}
\]
Elemento típico: $[(a,b)]$ --- par $\neq$ classe (\textbf{Lei~7}).
A condição de
equivalência é o \textbf{produto cruzado}: dois inteiros, um produto, nenhuma
divisão.

\medskip\noindent
\textbf{A estrutura dual} deste andar é realizada pela involução:
\[
\boxed{\;\iota(a,b)=(b,a).\;}
\]
Ela satisfaz $\iota^{2}=\mathrm{id}$ em todo $V$. No espaço de
representantes, o swap é total em $V=\mathbb{Z}^{2}$ e preserva $V^{\times}$,
mas não $V_{\mathrm{rat}}$ (porque $\iota(0,b)=(b,0)\notin V_{\mathrm{rat}}$).
O swap preserva $V^{\times}=\{(a,b):ab\neq0\}$ e desce ao quociente: se
$(a,b)\sim(c,d)$, isto é $ad=bc$, então $(b,a)\sim(d,c)$ porque $bc=ad$.
Portanto
\[
\bar\iota:V^{\times}/{\sim}\longrightarrow V^{\times}/{\sim},
\qquad\bar\iota([(a,b)])=[(b,a)]
\]
é bem definida. Resumindo:
\[
\boxed{
\iota:V^{\times}\to V^{\times}
\quad\Longrightarrow\quad
\bar\iota:V^{\times}/{\sim}\to V^{\times}/{\sim}
}
\]
Como $V^{\times}/{\sim}\cong\mathbb{Q}^{\times}$, esta aplicação é
precisamente $\bar\iota(x)=x^{-1}$. No domínio maior $V_{\mathrm{rat}}$, o
swap não é interno: $\iota(0,b)=(b,0)$ é a razão pela qual \textbf{não existe
inversão interna de $0$}. A extensão a $\mathbb{P}^{1}(\mathbb{Q})$ não faz parte da
construção de $\mathbb{Q}$; é a extensão projectiva que torna total a acção no
gume.

\medskip\noindent
\textbf{A Lei~$0$ via Möbius} (\code{algebrico def:lei0, thm:lei0,
thm:inversao-dual}). Na recta projectiva, o ponto é $[p:q]$ e a inversão é a
troca $S=\bigl(\begin{smallmatrix}0&1\\1&0\end{smallmatrix}\bigr)$,
$\det S=-1$:
\[
[p:q]\;\longmapsto\;[q:p],
\qquad
[0:1]\mapsto[1:0]=\infty,
\qquad
\boxed{\;0\dg=\infty.\;}
\]
Não se executa uma divisão: a inversão é a troca das coordenadas, sem teste e
sem ramo. A troca é a sua própria inversa ($S^{-1}=S$) porque
$\operatorname{adj}(S)=-S$ e $\det S=-1$, e os dois sinais cancelam:
$S^{-1}=\operatorname{adj}(S)/\det S = -S/(-1)=S$. Na extensão projectiva, os
dois pontos de gume são representados pelas classes $0=[0:1]$ e
$\infty=[1:0]$, que são trocadas pelo swap. As duas involuções trocam
exactamente os papéis de $\{0\}$ e de $\{\pm1\}$:
\[
\boxed{
\begin{array}{lll}
\text{em }\mathbb{P}^{1} & \text{fixos }\{+1,-1\}
  & \text{e o par }\{0,\infty\}\\[2pt]
\text{no trial} & \text{fixo }\{0\}
  & \text{e o par }\{-1,+1\}
\end{array}}
\]
O que num é ponto fixo, no outro é o par que se troca --- e é esta dualidade
que faz do zero o canónico (\code{algebrico prop:zero-dual-zero}).

\medskip\noindent\textbf{O gume: $\infty\mapsto-1$}
(\code{algebrico sec:lei0}). A Möbius $M=T_{-1}\circ S$, isto é $M(x)=1/x-1$,
satisfaz $M(\infty)=-1$ (com $\det M=-1$), e o seu ponto fixo resolve
$x^{2}+x-1=0$ --- para $m=1$, a razão áurea. Com a convenção de indexação e o operador $A_{m}$
definidos em \code{algebrico thm:cuspide}, a cadeia fica
$x_{0}=\infty$, $x_{-1}=0$, $x_{-2}=-1/m$.
Na interpretação herdada da torre, a operação que põe $\infty$ dentro da
aritmética aditiva é a mesma que gera o primeiro metal --- essa identificação
não é um teorema deste paper; é interface com o andar superior. Tudo isto é
\textbf{derivado}: $M$ não é primitiva; é a composição de $S$ (a Lei~$0$) com
a translação $T_{-1}$ (a Lei~$1$).

\medskip\noindent
As operações $J(a,b)=(-b,a)$ vivem em todo $V=\mathbb{Z}^{2}$.

\medskip\noindent
\textbf{Duas involuções distintas.} Há aqui duas involuções:
$\nu(a,b)=(-a,b)$, associada à reversibilidade aditiva (oposto), e
$\iota(a,b)=(b,a)$, associada à inversão multiplicativa (dual). A involução da
Lei~$1$ realiza a reversibilidade aditiva; o swap $\iota$ é a segunda
involução estrutural, responsável pela reversibilidade multiplicativa. Elas
não devem ser identificadas.

\medskip\noindent
\textbf{Quatro propriedades, em paralelo com o Algébrico}
(\code{algebrico def:espaco}).
\[
\begin{array}{ll@{\qquad}ll}
\textbf{(P1)} & \mathbb{Z}\hookrightarrow\mathbb{Q}\
  (\text{por }n\mapsto(n,1))
  & \textbf{(D1)} & \bar\iota([(a,b)])=[(b,a)]=[(a,b)]
    \iff a^{2}=b^{2}
    \iff [(a,b)]=\pm1,\; (a,b)\in V^{\times}\\[4pt]
\textbf{(P2)} & \text{nem todo racional é inteiro (e.g.\ }[(1,2)]=\tfrac12\notin\mathbb{Z}\text{)}
  & \textbf{(D2)} & [(a,b)]\otimes[(b,a)]
    =[(1,1)]\ \text{em }\mathbb{Q}^{\times}
\end{array}
\]
(P1) registra a inclusão do andar anterior:
$\mathbb{Z}\hookrightarrow\mathbb{Q}$, $n\mapsto[(n,1)]$.
(P2) registra que a extensão é própria: $\mathbb{Q}\neq\mathbb{Z}$. (D1) diz que os \textbf{pontos fixos da involução induzida} em
$\mathbb{Q}^{\times}$ são exactamente $\pm 1$: no espaço de representantes,
$\operatorname{Fix}(\iota|_{V})=\{(a,a):a\in\mathbb{Z}\}$;
mas \textbf{após o quociente},
$[(a,b)]=[(b,a)]\iff a^{2}=b^{2}\iff a=\pm b$ (porque $[(-1,1)]=[(1,-1)]$),
e os pontos fixos são
\[
\operatorname{Fix}(\iota\vert_{\mathbb{Q}^{\times}})
=\{x\in\mathbb{Q}^{\times}:x=x^{-1}\}
=\{\pm1\}.
\] (D2) diz que o produto pelo dual
\textbf{regressa ao elemento-base}: para $(a,b)\in V^{\times}$ (com
$ab\neq 0$), $(a,b)\otimes\iota(a,b)=(ab,ab)\sim(1,1)$,
logo o produto pelo dual regressa a $1$ \textbf{no quociente} --- o resultado
é literalmente diagonal antes do quociente. O coração do andar é:
\[
\boxed{
\iota(a,b)=(b,a)
\quad\Longrightarrow\quad
\bar\iota([(a,b)])=[(b,a)]=[(a,b)]^{-1}.
}
\]

\medskip\noindent
Compare com o andar de cima: em $V=\mathbb{Q}(m\sqrt D)$ a involução é
$s\mapsto s\dg=-1/s$ (a involução específica definida no andar quadrático,
\code{algebrico thm:fixo-dual}; não confundir com a conjugação de campo usual
$a+b\sqrt D\mapsto a-b\sqrt D$), com $s\,s\dg=-1$, e (D2) é
$\alpha\alpha^{*}\in\mathbb{Q}$. Aqui o swap dá
$(a,b)\otimes\iota(a,b)\sim(1,1)$, e (D2) é o produto pelo dual a regressar a
$1\in\mathbb{Z}$. \textbf{A semelhança é estrutural}: não se trata da mesma
involução, mas do mesmo \emph{molde} --- involução, ponto fixo, regresso ---
uma oitava abaixo.
\medskip\noindent
\textbf{Nota terminológica.} Neste paper, ``dual'' designa o espaço obtido
pela involução $\iota$, conforme a convenção da torre
(\code{analitico thm:tecidos}), e não o dual linear
$V^{*}=\operatorname{Hom}(V,\mathbb{Z})$.
\end{definicao}

\begin{proposicao}[$\sim$ é relação de equivalência]\label{prop:equiv}
Em $V_{\mathrm{rat}}$:
\begin{enumerate}\itemsep2pt
\item \textbf{Reflexiva:} $(a,b)\sim(a,b)$ pois $ab=ab$.
\item \textbf{Simétrica:} $ad=bc\Rightarrow cb=da$, logo $(c,d)\sim(a,b)$.
\item \textbf{Transitiva:} se $ad=bc$ e $cf=de$ com $d\neq 0$, então
      $af\cdot d=be\cdot d$, logo $af=be$ e $(a,b)\sim(e,f)$.
\end{enumerate}
\end{proposicao}

\begin{proof}
(1)--(2) imediatos. (3) só usa aritmética em $\mathbb{Z}$ --- sem invocar
$\mathbb{Q}$ antes do quociente estar definido. $\square$
\end{proof}

\begin{teorema}[Construção inteira de $\mathbb{Q}$]\label{thm:q}
Seja $V_{\mathrm{rat}}=\{(a,b)\in\mathbb{Z}^{2}:b\neq0\}$ com
$(a,b)\sim(c,d)\iff ad=bc$ (Prop.~\ref{prop:equiv}). Então o quociente
$V_{\mathrm{rat}}/{\sim}$ é um corpo; o mapa
\[
\boxed{\;
\Psi\colon V_{\mathrm{rat}}/{\sim}\;\overset{\cong}{\longrightarrow}\;\mathbb{Q},\qquad
\Psi([(a,b)])=a/b\;}
\]
é isomorfismo de corpos. Sob $\Psi$, as operações
\[
(a,b)\oplus(c,d)=(ad+bc,\;bd),\qquad
(a,b)\otimes(c,d)=(ac,\;bd)
\]
induzem a soma e o produto usuais, e todo elemento não nulo $[(a,b)]$ possui
inverso $[(a,b)]^{-1}=[(b,a)]$. Os elementos especiais são
\[
0=[(0,1)],\qquad 1=[(1,1)],\qquad -1=[(-1,1)],
\]
com $[(a,b)]\neq 0\iff a\neq 0$. Além disso,
\[
[(a,b)]=[(c,d)]\iff ad-bc=0,
\]
e, para $b,d>0$,
\[
[(a,b)]<[(c,d)]\iff ad-bc<0
\]
(\emph{ordem transportada} por $\Psi$ com denominadores positivos --- não ordem
canônica dos pares).
\end{teorema}

\begin{proof}[Prova]
(1) $\Psi$ bem-definida: Prop.~\ref{prop:dircr} (iii).
(2) Sobrejetiva: $q\in\mathbb{Q}$ --- representante $(p,q)$ com $q>0$.
(3) Injectiva: $\Psi([(a,b)])=\Psi([(c,d)])$ $\Rightarrow$ $a/b=c/d$
$\Rightarrow$ $ad=bc$ $\Rightarrow$ $[(a,b)]=[(c,d)]$.
(4) $\Psi$ preserva $\oplus$ e $\otimes$: Prop.~\ref{prop:bemdef} --- as operações
descem ao quociente.
(5) \textbf{Corpo no quociente.} Pelo Prop.~\ref{prop:bemdef}, $\oplus$ e $\otimes$ são
bem-definidas em $V_{\mathrm{rat}}/{\sim}$. Verificam-se os axiomas directamente:
\begin{itemize}\itemsep1pt
\item \emph{Adição:} $(0,1)$ neutro; $-(a,b)=(-a,b)$ oposto;
      associatividade e comutatividade herdam-se dos inteiros nos representantes.
\item \emph{Multiplicação:} $(1,1)$ neutro; $(a,b)\neq[(0,1)]\iff a\neq0$;
      inverso $[(b,a)]$ pois $(a,b)\otimes(b,a)\sim(1,1)$ quando $a\neq0$.
\item \emph{Distributividade:}
      $\Psi\bigl(x\otimes(y\oplus z)\bigr)=\Psi(x)\bigl(\Psi(y)+\Psi(z)\bigr)
      =\Psi(x)\Psi(y)+\Psi(x)\Psi(z)
      =\Psi(x\otimes y)+\Psi(x\otimes z)$ em $\mathbb{Q}$; pela injectividade de
      $\Psi$ (passo~(3)), $x\otimes(y\oplus z)=x\otimes y\oplus x\otimes z$.
\end{itemize}
Logo $(V_{\mathrm{rat}}/{\sim},\oplus,\otimes)$ \textbf{é} corpo; $\Psi$ isomorfismo
de corpos com $\mathbb{Q}$ usual.
(6) Ordem: comparar $ad$ e $bc$ com $b,d>0$. $\square$
\end{proof}

\begin{corolario}[Ordem no quociente]\label{cor:ordem}
Para representantes com $b,d>0$,
\[
x<y\iff\Psi(x)<\Psi(y)\iff\mathrm{Cruz}(x,y)<0.
\]
Ela \textbf{não} coincide com a ordem lexicográfica dos pares
$(a,b)\in\mathbb{Z}^{2}$.
\end{corolario}

\begin{observacao}[Molde do degrau anterior]
$\mathbb{N}\!\to\!\mathbb{Z}$ (\code{inteiros.tex}): $a+d=b+c$ (soma cruzada).
$\mathbb{Z}\!\to\!\mathbb{Q}$ (aqui): $ad=bc$ (produto cruzado). Analogia:
\[
\boxed{\text{soma cruzada}\longrightarrow\text{produto cruzado}}.
\]
\end{observacao}

\begin{teorema}[O swap é a inversão multiplicativa]\label{thm:swap-inverso}
A aplicação $\iota:V^{\times}\to V^{\times}$, $\iota(a,b)=(b,a)$, preserva
$\sim$, induz uma involução
$\bar\iota:\mathbb{Q}^{\times}\to\mathbb{Q}^{\times}$, e sob o isomorfismo
$[(a,b)]\mapsto a/b$ satisfaz $\bar\iota(x)=x^{-1}$, com
$x\,\bar\iota(x)=1$ e $\operatorname{Fix}(\bar\iota)=\{\pm1\}$.
\end{teorema}

\noindent O objecto novo deste degrau não é uma nova operação sobre números: é
a representação que torna a inversão multiplicativa uma involução inteira no
espaço de pares:
\[
\boxed{
\begin{array}{c}
\text{pares inteiros}\\
\downarrow\\
\text{produto cruzado}\\
\downarrow\\
\text{classe de equivalência}\\
\downarrow\\
\mathbb{Q}\\
\downarrow\\
\text{swap}\\
\downarrow\\
x\mapsto x^{-1}
\end{array}}
\]

\noindent A realização computacional em \code{Qz} implementa este teorema sob
um envelope finito explicitamente medido.

%======================================================================
\section{As operações no par --- sem sair de inteiros}\label{sec:ops}

\begin{definicao}[aritmética dos pares]\label{def:ops}
Sejam $(a,b)$ e $(c,d)$ representantes com $b,d\neq 0$.
\[
\boxed{\;
\begin{array}{rll}
\textbf{soma:}     & (a,b)\oplus(c,d) &= (ad+bc,\; bd)\\[4pt]
\textbf{produto:}  & (a,b)\otimes(c,d) &= (ac,\; bd)\\[4pt]
\textbf{oposto:}   & -(a,b) &= (-a,\; b)\\[4pt]
\textbf{inverso:}  & (a,b)^{-1} &= (b,\; a) \qquad(a\neq 0)
\end{array}\;}
\]
\end{definicao}

\noindent Cada uma produz um \textbf{novo par de inteiros}. O inverso é o
\textbf{swap} $\iota$ --- a mesma involução da Def.~\ref{def:espaco}.
No quociente, escrevemos
$[(a,b)]\otimes[(c,d)]:=[(ac,bd)]$, e analogamente para as demais operações.

\begin{proposicao}[boa definição]\label{prop:bemdef}
As quatro operações são \textbf{bem definidas} sobre classes: se
$(a,b)\sim(a',b')$ e $(c,d)\sim(c',d')$, então os resultados são equivalentes.
Como $\operatorname{eval}$ tem precisamente $\sim$ como relação de fibras
($ab'=a'b \Leftrightarrow a/b=a'/b'$ para $b,b'\neq 0$), representantes
equivalentes produzem o mesmo racional; portanto as operações descendem ao
quociente. Para o produto, por exemplo: $ab'=a'b$ e $cd'=c'd$ implicam
$(ac)(b'd')=(a'c')(bd)$. Para o inverso: $ab'=a'b$ implica $ba'=b'a$,
portanto $(b,a)\sim(b',a')$. $\square$
\end{proposicao}

%======================================================================
\section{A transformada e o ciclo de operação/inversão}\label{sec:conv}

No Algébrico (\code{algebrico thm:universal}), o ciclo universal tem quatro
passos: \textsc{descodifica}, \textsc{opera}, \textsc{inverte},
\textsc{codifica}. O passo \textsc{opera} é a \textbf{convolução} e o passo
\textsc{inverte} é a \textbf{deconvolução}. Neste andar, tomamos a avaliação $\operatorname{eval}$ como a realização da
transformada; sob essa realização, as identidades de soma e produto descendem
imediatamente.

\begin{definicao}[a transformada neste degrau]\label{def:transf}
No Algébrico, a transformada universal é a avaliação nas folhas
$\sigma,\sigma\dg$ (\code{algebrico thm:espectro}). Neste degrau, a
``folha'' é a avaliação $\operatorname{eval}(a,b)=a/b$, que é a projecção
$V_{\mathrm{rat}}\to\mathbb{Q}$. Sob a avaliação, o produto de pares é
levado ao produto dos valores:
\[
\operatorname{eval}\bigl((a,b)\otimes(c,d)\bigr)
= \frac{ac}{bd}
= \frac{a}{b}\cdot\frac{c}{d}
= \operatorname{eval}(a,b)\cdot\operatorname{eval}(c,d).
\]
A convolução, após a avaliação na única folha, reduz-se ao produto ordinário
de $\mathbb{Q}$. A soma:
\[
\operatorname{eval}\bigl((a,b)\oplus(c,d)\bigr)
= \frac{ad+bc}{bd}
= \frac{a}{b}+\frac{c}{d}
= \operatorname{eval}(a,b)+\operatorname{eval}(c,d).
\]
A avaliação é compatível com $\oplus$ e $\otimes$ e induz no quociente o
isomorfismo de anéis
$\overline{\operatorname{eval}}:V_{\mathrm{rat}}/{\sim}\;\xrightarrow{\sim}\;\mathbb{Q}$,
$[(a,b)]\longmapsto a/b$. A palavra ``transformada'' é terminologia interna da
série; matematicamente, $\operatorname{eval}$ é a aplicação canónica de
representantes para o corpo de fracções (não se trata de uma transformada
integral ou espectral).
\end{definicao}

\noindent\textbf{Nota terminológica.} Nesta seção, ``convolução'' e
``deconvolução'' são nomes internos do ciclo estrutural da série; nenhuma
propriedade da convolução analítica clássica é usada ou reivindicada.

\begin{proposicao}[produto e inversão --- o ciclo reversível]\label{prop:conv}
Neste degrau, o par estrutural correspondente ao ciclo operação/inversão
realiza-se pelo produto de representantes $(a,b)\otimes(c,d) = (ac,\;bd)$
e pela involução swap $(a,b)^{-1} = (b,a)$, $a\neq 0$, que implementa o
inverso multiplicativo. O par fecha:
$(a,b)\otimes\iota(a,b)=(ab,ab)\sim(1,1)=1$.
O produto desce sempre; a inversão desce \textbf{sob condição}
($a\neq 0$) --- a mesma estrutura do
\code{algebrico thm:decon-andar}.
No vocabulário interno da série (\code{analitico thm:bit}), este par
corresponde ao ciclo convolução (gato) / deconvolução (esquilo).

A Möbius não é importada como primitiva: é \textbf{derivada} do dual
(\code{algebrico thm:inversao-dual}). Para qualquer matriz invertível
$2\times2$, $M^{-1}=\operatorname{adj}(M)/\det M$; e para a troca
$S=\bigl(\begin{smallmatrix}0&1\\1&0\end{smallmatrix}\bigr)$, com
$|\det S|=1$, a inversa fica em $\mathbb{Z}$. É por isso que o swap inteiro
nos pares realiza a inversão multiplicativa sem sair dos inteiros. $\square$
\end{proposicao}

%======================================================================
\section{Dir e Cruz: componentes simétrica e antissimétrica da leitura}\label{sec:operacao}

Da transformada (Def.~\ref{def:transf}) derivam-se as duas leituras: Dir e
Cruz são os \textbf{numeradores inteiros das leituras de soma e diferença
quando colocadas sobre o denominador comum $bd$}. A leitura bilinear $B$ e a
troca de seus argumentos produzem a forma transposta $B^{\tau}$; os dois
acoplamentos resultantes são $ad$ e $bc$.

\begin{definicao}[Dir e Cruz como decomposição da
transformada]\label{def:dircr}
Escolha-se a forma bilinear $B((a,b),(c,d))=ad$, compatível com o primeiro
acoplamento da soma em denominador comum; a troca dos argumentos induz a forma
transposta $B^{\tau}(x,y)=B(y,x)$:
\[
B\bigl((a,b),(c,d)\bigr) = ad,\qquad
B^{\tau}\bigl((a,b),(c,d)\bigr) = bc.
\]
A troca dos argumentos da forma bilinear produz a forma transposta $B^{\tau}$,
e os dois acoplamentos $ad$ e $bc$ são precisamente os termos obtidos dessa
dupla leitura. $B$ fornece o primeiro dos dois acoplamentos inteiros que
compõem o numerador da soma no denominador comum $bd$.
As partes simétrica e antissimétrica são:
\[
\boxed{\;
\operatorname{Dir} = B + B^{\tau} = ad+bc,
\qquad
\operatorname{Cruz} = B - B^{\tau} = ad-bc.
\;}
\]
Após extensão dos escalares a $\mathbb{Q}$,
$P_{\pm}:\operatorname{Bil}(V,V;\mathbb{Q})\to
\operatorname{Bil}(V,V;\mathbb{Q})$, $P_{\pm}(B)=\frac{1}{2}(B\pm B^{\tau})$,
são os projectores sobre as partes simétrica e antissimétrica do espaço de
formas bilineares. Sobre $\mathbb{Z}$, Dir e Cruz são os numeradores inteiros
correspondentes, $\operatorname{Dir},\operatorname{Cruz}\in
\operatorname{Bil}(\mathbb{Z}^{2},\mathbb{Z})$:
$\operatorname{Dir}=2P_{+}(B)$, $\operatorname{Cruz}=2P_{-}(B)$.
Os projectores $P_{\pm}$ são introduzidos apenas como interpretação racional
da decomposição; Dir e Cruz permanecem definidos integralmente.

\medskip\noindent\textbf{A derivação pela transformada.} No Algébrico
(\code{analitico cor:dir-cruz-folhas}), Dir é o que a transformada vê --- o produto
folha a folha --- e Cruz é o que sobra quando as folhas coincidem. Neste
degrau, com uma única folha $\operatorname{eval}$:
\begin{itemize}
\item $\operatorname{Dir}((a,b),(c,d)) = ad+bc$: é o numerador de
  $\operatorname{eval}(a,b)+\operatorname{eval}(c,d)$ no denominador comum
  $bd$. \textbf{Dir é a soma vista pela transformada.}
\item $\operatorname{Cruz}((a,b),(c,d)) = ad-bc$: é o numerador de
  $\operatorname{eval}(a,b)-\operatorname{eval}(c,d)$ no denominador comum
  $bd$. \textbf{Cruz é a diferença vista pela transformada.}
\end{itemize}
As duas leituras são construídas a partir da forma bilinear e da sua
transposição induzida pelo swap;
não são axiomas adicionais.
\end{definicao}

\begin{proposicao}[as propriedades são as do Algébrico]\label{prop:dircr}
\emph{(i)} Dir é simétrico: $ad+bc=cb+da$. \emph{(ii)} Cruz é
antissimétrico: $ad-bc=-(cb-da)$. \emph{(iii)} $\operatorname{Cruz}=0 \iff
ad=bc \iff (a,b)\sim(c,d)$: a igualdade de racionais \textbf{é} o cruzado
nulo. \emph{(iv)} A \textbf{ordem} (com $b,d>0$) é
$a/b < c/d \iff ad-bc<0$. $\square$
\end{proposicao}

\noindent O produto cruzado não é apenas uma implementação da equivalência:
\textbf{é a própria testemunha inteira da equivalência}. A relação
$(a,b)\sim(c,d)\iff ad-bc=0$ é o núcleo do paper.

\begin{observacao}[por que uma folha e não duas]
No andar quadrático, $\mathbb{Q}(m\sqrt D)$ tem \textbf{duas} folhas
$\sigma,\sigma\dg$, e a dourada discreta (\code{algebrico
thm:dourada-discreta}) opera com ambas. Neste degrau, para cada representante $(a,b)$ com $b\neq0$, a
relação $bx-a=0$ é linear sobre $\mathbb{Q}$ e determina a única raiz
$x=a/b$. Neste modelo, a realização da transformada possui uma única folha
racional, dada pela avaliação $a/b$, e o par
convolução\,/\,deconvolução é exactamente o produto\,/\,divisão de
racionais. No modelo quadrático adotado no andar seguinte, essa passagem
manifesta-se como a introdução de duas folhas $\sigma,\sigma\dg$
(\code{algebrico thm:fixo-dual}).
\end{observacao}

%======================================================================
\section{O gume: a divisão por zero}\label{sec:gume}

O gume deste andar é a \textbf{ausência de fibra} na divisão por zero:
$0\cdot x = 0$ para todo $x$, logo nenhum $x$ cumpre $0\cdot x=c$ com
$c\neq 0$. O swap $\iota(0,b)=(b,0)$ \textbf{sai do espaço racional}
($b\neq 0$ mas $0$ na segunda posição): é exactamente a Lei~$0$ a levar $0$
a $\infty$ em $\mathbb{P}^{1}$. Por isso, o swap realiza uma involução de
$\mathbb{Q}^{\times}$, não de $\mathbb{Q}$, pois $0$ não possui inverso
(\code{lib/racionais.h}, \code{qz\_inverso}):
\[
\boxed{
\begin{array}{c}
\mathbb{Q}^{\times}
\overset{\iota}{\longrightarrow}
\mathbb{Q}^{\times}
\qquad\text{(swap nos invertíveis)},\\[4pt]
\mathbb{P}^{1}(\mathbb{Q})
\overset{\iota}{\longrightarrow}
\mathbb{P}^{1}(\mathbb{Q})
\quad\text{(swap total)}.
\end{array}}
\]

%======================================================================
\section{A redução: o fator comum removido}\label{sec:reduz}

O par $(a,b)$ não se guarda como calhou: \textbf{reduz-se} pelo $\gcd$, que
remove a redundância de representação, com a convenção canónica $b>0$ e
$\gcd(\lvert a\rvert,b)=1$:
\[
(a,b)\mapsto\left(\frac{a}{g},\;\frac{b}{g}\right),\qquad g=\gcd(a,b),
\quad\text{com sinal ajustado para }b>0.
\]
No produto $(a,b)\otimes(c,d)$, cancela-se em \textbf{cruz}:
$g_{1}=\gcd(a,d)$, $g_{2}=\gcd(c,b)$, e substitui-se
$a\leftarrow a/g_{1}$, $d\leftarrow d/g_{1}$,
$c\leftarrow c/g_{2}$, $b\leftarrow b/g_{2}$,
\emph{antes} de multiplicar. Isto evita
crescimento artificial devido à redundância de representação; a permanência
dentro de um envelope de $16$~bits é uma propriedade adicional,
verificada pelos testes da cadeia.

%======================================================================
\section{As oito leis neste andar}\label{sec:leis}

\noindent Onde cada lei vive na cadeia --- tabela comum aos três papers
--- a mesma nos três:

\begin{center}\footnotesize
\begin{tabular}{@{}cp{2.75cm}p{2.75cm}p{2.75cm}p{2.75cm}@{}}
\toprule
lei & $k{=}-1$: $F_1\!\to\!F_8$ & $k{=}0$: $\mathbb{N}\!\to\!\mathbb{Z}$ & $k{=}1$: $\mathbb{Z}\!\to\!\mathbb{Q}$ & $k{=}2$: $\mathbb{Q}\!\to\!\mathbb{R}$ \\
\midrule
\textbf{0} divisão do zero
 & $x\oplus x=0$; só o $0$ sem inversa \hfill\textsc{expl.}
 & subtracção parcial; folha $[(1,0)]\oplus[(0,1)]=0$ \hfill\textsc{expl.}
 & $0^{-1}$ não existe; $0\!\leftrightarrow\!\infty$ \hfill\textsc{mista}
 & o \textbf{buraco}: corte sem racional em cima \hfill\textsc{expl.} \\
\textbf{1} dual ($1\dg=-1$)
 & \emph{degenera}: $-x=x$ \hfill\textsc{expl.}
 & $\iota$ no par; $\nu(z)=-z$ \hfill\textsc{expl.}
 & $\nu(a,b)=(-a,b)$ \hfill\textsc{expl.}
 & $\nu(x)=-x$; do andar: $(A,B)\!\mapsto\!(B,A)$ \hfill\textsc{expl.} \\
\textbf{2} bidual
 & $x^{**}=x$; automorfismo \hfill\textsc{expl.}
 & Dir/Cruz; $R^{4}=\mathrm{id}$ \hfill\textsc{mista}
 & $V\cong V^{**}$ canónico \hfill\textsc{canón.}
 & $B=\mathbb{Q}\setminus A$: complementar $2\times$ \hfill\textsc{expl.} \\
\textbf{3} trial $\{-1,0,+1\}$
 & \emph{degenera}: sem cúspide \hfill\textsc{expl.}
 & cúspide $\tau=0$; $+0=-0$ \hfill\textsc{expl.}
 & $\operatorname{sgn}(a/b)$ \hfill\textsc{expl.}
 & $\tau=0$ é \emph{vazio} $\iff$ irracional \hfill\textsc{expl.} \\
\textbf{4} tetral $T+T^{*}$
 & $F_{2w}=F_w\oplus\sigma F_w$ \hfill\textsc{expl.}
 & $\mathbb{N}^{2}/{\sim}$; classe $[(a,b)]$ \hfill\textsc{expl.}
 & recebida (\code{thm:juntas}) \hfill\textsc{receb.}
 & $(A,B)$, e quocientar aberto/fechado \hfill\textsc{expl.} \\
\textbf{5} pental $x^{2}=-1$
 & \emph{degenera}: uma raiz só \hfill\textsc{expl.}
 & $G_{\mathrm{geom}}\bmod2$ \hfill\textsc{receb.}
 & $J(a,b)=(-b,a)$, $J^{2}=-I$ \hfill\textsc{expl.}
 & \emph{não existe}: ordenado \hfill\textsc{expl.} \\
\textbf{6} hexal (interface)
 & base ortonormal: coordenada $=$ ler o bit \hfill\textsc{receb.}
 & $\otimes\leftrightarrow(a-b)(c-d)$ \hfill\textsc{receb.}
 & interface \hfill\textsc{receb.}
 & densidade de $\mathbb{Q}$ em $\mathbb{R}$ \hfill\textsc{receb.} \\
\textbf{7} octonião dual
 & encaixe por prefixos \hfill\textsc{expl.}
 & $q\parallel p$; par $\neq$ classe; $\Word_8$ \hfill\textsc{expl.}
 & junta dual \hfill\textsc{receb.}
 & três realizações ligadas sem fundidas \hfill\textsc{expl.} \\
\bottomrule
\end{tabular}
\end{center}
\noindent{\footnotesize \textsc{expl.} realizada neste degrau ·
\textsc{receb.} interface herdada do centro, não derivada aqui ·
\textsc{mista}/\textsc{canón.} ver o paper do degrau.
\textbf{Três leis degeneram em $k=-1$} (1, 3, 5) --- e é isso que permite a
torre binária dobrar sem excepção. A Lei~5 falta nos dois extremos por razões
\emph{opostas}: em $k=-1$ por característica~2, em $k=2$ por ordem.
\textbf{Nenhum degrau realiza as oito}; a coluna diz onde cada uma vive.\par}

\noindent E, neste andar, com o domínio de cada realização.
As leis não são redefinidas axiomaticamente: este paper \textbf{realiza}, no
espaço de pares, as operações herdadas da construção universal
(\code{analitico sub:oitoleis}). Quando a operação é explicitada, a preservação
do domínio inteiro é imediata; quando depende de estrutura adicional, a
referência é ao teorema correspondente do \emph{Corpo Algébrico}.

\begin{center}\small
\begin{tabular}{@{}clll@{}}
\toprule
lei & realização no par & domínio & estatuto \\
\midrule
\textbf{0} & $0=(+1)\oplus(-1)$ em $\mathbb{Q}$; $0\leftrightarrow\infty$ em $\mathbb{P}^{1}$ &
  $\mathbb{Q}$, $\mathbb{P}^{1}(\mathbb{Q})$ & mista \\
\textit{dual do andar} & $\iota(a,b)=(b,a)$, $\iota^{2}=I$ (involução multiplicativa) &
  $\mathbb{Q}^{\times}$ & estrutura do degrau \\
\textbf{1} & $\nu(a,b)=(-a,b)$, $\nu^{2}=I$ & $\mathbb{Z}^{2}$ & explícita \\
\textbf{2} & $V\xrightarrow{\sim} V^{**}$ (dual linear canónico) & $\mathbb{Z}^{2}$ & canónica \\
\textbf{3} & $\operatorname{sgn}(a/b)\in\{-1,0,+1\}$ & $\mathbb{Q}$ & explícita \\
\textbf{4} & $T+T^{*}$ (tecidos) & --- & \code{algebrico thm:juntas} \\
\textbf{5} & $J(a,b)=(-b,a)$, $J^{2}=-I$ & $\mathbb{Z}^{2}$ & explícita \\
\textbf{6} & operação de interface (hexal) & --- & herdada \\
\textbf{7} & junta dual (octonião) & --- & herdada \\
\bottomrule
\end{tabular}
\end{center}

\noindent As Leis $1$, $3$ e $5$ possuem \textbf{realização interna imediata}
em $\mathbb{Z}^{2}$; a Lei~$0$ possui realização aditiva interna
($0=(+1)\oplus(-1)$) e realização do gume apenas após a extensão projectiva; a Lei~$2$ é canónica: $V=\mathbb{Z}^{2}\overset{\sim}{\longrightarrow}V^{**}$
pelo mapa $v\mapsto(\phi\mapsto\phi(v))$, que é isomorfismo porque $V$ é livre
de posto finito, e não depende do swap. (A palavra ``dual'' na Lei~$2$ refere-se
ao dual linear $V^{*}=\operatorname{Hom}(V,\mathbb{Z})$, não ao dual da torre
$\iota$, nem ao dual quadrático $s^{\dagger}=-1/s$ do Corpo Algébrico; as três
noções não devem ser confundidas.) As Leis $4$, $6$ e $7$ pertencem à estrutura global recebida da torre; este
paper não fornece uma nova realização independente delas --- são registadas
para explicitar a posição do degrau na hierarquia. A identidade $J^{4}=I$ é consequência da
Lei~$5$: $J^{2}=-I$ é a realização em $\operatorname{End}(\mathbb{Z}^{2})$ da
relação $x^{2}=-1$, portanto $J^{4}=I$.

As realizações são inteiras em $\mathbb{Z}^{2}$; quando o resultado deve
permanecer em $\mathbb{Q}$, a Lei~$0$ aplica-se em $\mathbb{Q}^{\times}$, e a
sua extensão natural é $\mathbb{P}^{1}(\mathbb{Q})$.

%======================================================================
\section{A ligação com o andar quadrático}\label{sec:zsigma}

No andar de cima, o par $(a,b)$ é \textbf{coordenada linear na base}:
$\alpha=a+b\sigma$ em $\mathbb{Q}(\sigma)$. Neste andar, o par $(a,b)$ é
\textbf{coordenada homogénea}: $[a:b]$ com $b\neq 0$ dá $\mathbb{Q}$
(a parte afim de $\mathbb{P}^{1}(\mathbb{Q})$); incluindo $[1:0]=\infty$
obtém-se $\mathbb{P}^{1}(\mathbb{Q})$.
São dois papéis diferentes do mesmo dado $(a,b)$:
\[
\boxed{\;
\begin{array}{ll}
\mathbb{Q}: & [a:b],\ b\neq 0\\[3pt]
\mathbb{P}^{1}(\mathbb{Q}): & [a:b]\ \text{incluindo }[1:0]=\infty\\[3pt]
\mathbb{Q}(\sigma): & a+b\sigma \quad\text{(coordenada linear)}
\end{array}\;}
\]
No andar racional, o swap realiza a inversão multiplicativa e normaliza o
produto dual para $1$. No andar quadrático, a involução específica
$s\mapsto s\dg=-1/s$ satisfaz $s\,s\dg=-1$; ela \textbf{não} é a conjugação
de campo $a+b\sqrt{D}\mapsto a-b\sqrt{D}$, embora ambas sejam involuções:
\[
\begin{array}{lll}
\text{conjugação} & a+b\sqrt{D}\mapsto a-b\sqrt{D} & N(s)\in\mathbb{Q},
\quad\operatorname{Fix}=\mathbb{Q}\\[2pt]
\text{dobra }\dagger & s\mapsto -1/s & s\,s\dg=-1,
\quad\operatorname{Fix}=\{s:s^{2}=-1\}
\end{array}
\]
A semelhança é \emph{estrutural} --- não se trata da mesma involução.

\subsection*{O ponto fixo e a descida}

Na parametrização específica do \emph{Corpo Algébrico}
(\code{algebrico thm:fixo-dual}), o ponto fixo de $g(x)=m+1/x$ é o vector
próprio de $A_m$, raiz de $x^{2}=mx+1$. O par dual $\sigma,\sigma\dg$
satisfaz $\sigma+\sigma\dg=m$ (traço, inteiro) e $\sigma\,\sigma\dg=-1$
(norma, constante). Na parametrização adotada, o inteiro associado ao ponto
fixo é o seu traço, e a classe de equivalência é realizada, no modelo
quadrático adotado, pela fibra do traço
(\code{algebrico thm:descida}):
\[
\mathbb{Z}\;\cong\;\{\text{pontos fixos}\}/{\sim},
\qquad
\sigma\sim\tau\iff\operatorname{tr}\sigma=\operatorname{tr}\tau.
\]
Na parametrização específica adoptada no \emph{Corpo Algébrico}, em que
$\operatorname{tr}(\sigma_m)=m$, o racional aparece como \textbf{razão de
traços}: $a/b = \operatorname{tr}(\sigma_a) / \operatorname{tr}(\sigma_b)$; e a relação $(a,b)\sim(c,d)\iff ad=bc$ é
\emph{recuperada} dos traços cruzados. Esta identificação não é uma construção
de $\operatorname{Frac}(\mathbb{Z})$; é uma realização particular da interface
com o andar quadrático já construída no \emph{Corpo Algébrico}.
Dentro de cada corpo $\mathbb{Q}(m\sqrt D)$, a cópia de $\mathbb{Q}$ é o
auto-espaço $+1$ da \textbf{conjugação de campo}
$a+b\sqrt{D}\mapsto a-b\sqrt{D}$, não da involução específica
$s\dg=-1/s$ (cujo ponto fixo satisfaria $s^{2}=-1$)
(\code{algebrico thm:descida}(4)).

\subsection*{A posição na tríade: Hurwitz, Gentil e Lebesgue}

\noindent A identificação que se segue com a tríade é uma interpretação
herdada da arquitectura global; não é usada na prova da construção de
$\mathbb{Q}$.

No vocabulário da série, o teorema central (\code{analitico thm:central}) diz
que \textbf{contar} (Hurwitz, o lado discreto --- a norma
$N(x)=\sum x_i^{2}$) e \textbf{integrar} (Gentil, o lado contínuo --- a fusão
homogénea) são duais
pela medida. Se $f:[a,b]\to[f(a),f(b)]$ é uma bijecção monotónica contínua e
as integrais envolvidas existem, a soma reversível de Lebesgue dá:
\[
\int_a^b f \;+\; \int_{f(a)}^{f(b)} f^{-1} \;=\; b\,f(b)-a\,f(a).
\]
Neste degrau, a posição é precisa:
\begin{itemize}
\item O par $(a,b)$ e o produto cruzado $ad-bc$ vivem no lado
  \textbf{discreto}: são contagem inteira, sem raiz, sem limite --- Hurwitz.
\item A completação $\mathbb{Q}\to\mathbb{R}$ pelo corte é o salto para o
  lado \textbf{contínuo}: o ponto fixo que falta ao andar racional
  (\code{algebrico thm:corte-fixo}) é alcançado pela soma reversível, não por
  uma raiz --- Gentil.
\item No vocabulário da série, a preservação da ordem pelo produto cruzado
  (sob $b,d>0$) é a interface que será interpretada no andar contínuo como
  parte do transporte associado à medida (Lebesgue).
\end{itemize}
No vocabulário da tríade adoptado nesta série, o degrau
$\mathbb{Z}\to\mathbb{Q}$ pertence ao lado discreto (Hurwitz): \textbf{tudo é
contagem}. O salto seguinte $\mathbb{Q}\to\mathbb{R}$ é onde Gentil entra, e
a reversibilidade da soma é o mecanismo que liga os dois sem avaliar uma raiz.

%======================================================================
\section{O zero como cúspide e ponto fixo}\label{sec:cuspide}

\noindent\textbf{Nota.} As identificações desta secção são estrutura herdada
da geometria do \emph{Corpo Algébrico}; não são necessárias para a construção
de $\mathbb{Q}$ apresentada no Teorema~\ref{thm:q} e são usadas aqui apenas
para posicionamento na torre.

\medskip\noindent
No \emph{Corpo Algébrico} (\code{algebrico thm:cuspide}), a cúspide é o lugar
onde o discriminante se anula e as duas folhas da transformada universal
colidem: $\operatorname{disc}M=\operatorname{tr}^{2}M-4\det M=0$.
No formalismo geométrico herdado, os racionais aparecem como dados associados
à cúspide, e os cortes irracionais são o complemento.

\medskip\noindent\textbf{A inversão é derivada, não primitiva} --- cf.\
Prop.~\ref{prop:conv}.

\medskip\noindent\textbf{O zero é o ponto fixo da reflexão.} A tríade
$\tau\in\{-1,0,+1\}$ indexa os regimes (elíptico, parabólico, hiperbólico). A
reflexão $\tau\mapsto-\tau$ é uma involução, e o seu único ponto fixo em
$\mathbb{Z}$ é $\tau=0$:
\[
\boxed{\;\tau\dg=\tau\iff -\tau=\tau\iff 2\tau=0\iff\tau=0.\;}
\]
É por isso que o zero é canónico: não se escolhe a origem --- \emph{quem se
fixa é a origem}. No nível dos pares, $0=[(0,1)]$ é a classe cujo swap
$(1,0)\notin V_{\mathrm{rat}}$ escapa para a extensão projectiva; a
Lei~$0$ é a realização projectiva deste gume.

\medskip\noindent\textbf{Observação --- interpretação herdada da torre.}
\[
\begin{array}{clll}
\tau=-1 & \text{elíptico} & \textbf{topologia} & \text{fecha: ordem finita, tudo volta}\\[2pt]
\tau=\phantom{-}0 & \text{cúspide} & \textbf{álgebra} & \text{ponto fixo: só }\pm I\text{ fecham}\\[2pt]
\tau=+1 & \text{hiperbólico} & \textbf{análise} & \text{não fecha: é o corte}
\end{array}
\]
A reflexão troca topologia com análise e fixa a álgebra --- é a dualidade
discreto\,$\leftrightarrow$\,contínuo com o fundamento quieto. A
identificação entre esses dois papéis do zero --- $0=[(0,1)]$ como neutro
aditivo de $\mathbb{Q}$ e $\tau=0$ como ponto fixo da reflexão --- é uma
\textbf{compatibilidade herdada} da estrutura global; não é necessária para a
construção do corpo $\mathbb{Q}$. Na interpretação herdada da torre, este
degrau está localizado na cúspide ($\tau=0$); a passagem ao corte ($\tau=+1$)
é o degrau seguinte.

%======================================================================
\section{O degrau seguinte: a completação}\label{sec:qstar}

O degrau seguinte da torre é
$\mathbb{Q}\longrightarrow\mathbb{R}$.
Neste formalismo, $\mathbb{R}$ é obtido pela adjunção dos pontos determinados
pelos cortes de $\mathbb{Q}$. Na linguagem da torre
(\code{analitico thm:tecidos}):
\[
T_{k+1}=T_{k}+\mathcal{C}(T_{k}),
\]
onde $\mathcal{C}(\mathbb{Q})$ é o conjunto de cortes que $\mathbb{Q}$ não
contém. É neste degrau que surge a completude necessária à passagem do domínio
enumerável para a estrutura contínua de $\mathbb{R}$.

\subsection*{Três caminhos para o mesmo ponto}

O \code{lib/reais.h} fornece três realizações do mesmo valor real $\sqrt{a}$;
a cadeia de testes verifica que elas concordam. A equivalência entre essas
realizações é uma propriedade da interface com o andar seguinte; não é
necessária para o Teorema~\ref{thm:q}.

\begin{center}\small
\begin{tabular}{@{}lll@{}}
\toprule
caminho & construção & tipo \\
\midrule
\textbf{1. o corte} & $A=\{q\in\mathbb{Q}:q\leq 0 \text{ ou } q^2<a\}$ --- decisão inteira &
\code{Corte \{a, n\}} \\
\textbf{2. o ponto fixo} & $x\mapsto(a+bx)/(x+b)$, Möbius --- cada iterado é racional &
\code{rz\_passo} \\
\textbf{3. a fração contínua} & periódica por Lagrange & \code{lado} \\
\bottomrule
\end{tabular}
\end{center}

\noindent As etapas finitas das três realizações usam a aritmética de pares
deste paper; o objecto limite pertence à completude.

\subsection*{O que a completação acrescenta à migração}

\begin{enumerate}
  \item As decisões finitas que constroem o corte não exigem um tipo aritmético
        novo: são realizadas sobre pares de inteiros. O objecto completo
        determinado pelo corte pertence ao novo domínio $\mathbb{R}$.
  \item A bisseção (\code{rz\_encaixota}) e o Möbius (\code{rz\_passo}) operam
        sobre \code{Qz} --- o par deste paper.
  \item Neste formalismo, o salto $\mathbb{Q}\to\mathbb{R}$ é o primeiro desta
        cadeia em que a completude passa a ser uma propriedade estrutural
        necessária.
\end{enumerate}

%======================================================================
\section{Irracionalidade de $\sqrt{2}$}\label{sec:sqrt2}

O buraco que $\mathbb{Q}$ não contém é visível pelo \textbf{corte}.
Defina
\[
A = \{r\in\mathbb{Q}: r\leq 0 \text{ ou } r^2 < 2\}.
\]
$A$ é um corte de Dedekind: não vazio, fechado para baixo, sem maior elemento.
Para ver que $A$ não tem máximo: se $r\in A$ com $r>0$ e $r^{2}<2$, escolha
racional $0<h<\min\{1,\,(2-r^{2})/(2r+1)\}$; como $h<1$ implica $h^{2}<h$,
$(r+h)^{2}=r^{2}+2rh+h^{2}<r^{2}+(2r+1)h<2$, logo $r+h\in A$ e $r+h>r$.
O real determinado por esse corte não pertence a $\mathbb{Q}$: se $s=p/q$ com $\gcd(p,q)=1$ e $s>0$,
então $s^{2}=2$ implica $p^{2}=2q^{2}$, que obriga $p$ par; escrevendo $p=2k$,
$4k^{2}=2q^{2}$ obriga $q$ par --- contradição com $\gcd(p,q)=1$.

\noindent Logo o corte $A$ define um real que não é racional. A decisão
``$r^{2}<2$?'' é inteira: $(p,q)\mapsto p^{2}$ vs $2q^{2}$, produto cruzado
de pares --- \textbf{nenhum limiar, nenhuma tolerância}.

\begin{observacao}[o MMC não pertence ao núcleo]
A construção de $\mathbb{Q}$ por pares de inteiros não requer uma unidade
comum de denominadores. O MMC pertence somente ao caminho de representação que
escolhe uma escala comum para vários racionais. No núcleo algébrico, soma,
produto, inversão e comparação operam directamente sobre pares, com
cancelamento por $\gcd$ e produto cruzado:
\[
\text{MMC é uma normalização opcional de representação;}\qquad
\gcd + \text{Cruz são suficientes para a construção.}
\]
\end{observacao}

%======================================================================
\section{Resumo para a migração}\label{sec:migra}

\begin{enumerate}
  \item O racional \textbf{é a classe} de um par ordenado de inteiros $(a,b)$,
        $b\neq 0$.
  \item As operações fazem-se por aritmética de inteiros: soma, produto, oposto,
        inverso (swap).
  \item A ordem é por produto cruzado: $ad$ vs $bc$, com $b,d>0$.
  \item Reduzir pelo $\gcd$ remove a redundância de representação.
  \item As leis $0$--$3$ e $5$ têm realizações inteiras explícitas em
        $\mathbb{Z}^{2}$.
  \item O tipo computacional é o \code{Qz \{int p, q;\}} de
        \code{lib/racionais.h}, com comparação por \code{d32\_cmp\_prod}.
  \item Para descer a $16$ bits, a representação passa a ser
        \code{Q16 \{int16 p, q;\}}, mantendo a mesma estrutura de pares. Os
        produtos intermediários continuam a ser realizados em largura superior
        ($32$ bits) e reduzidos antes de voltar ao domínio de $16$ bits. A
        validade da migração depende de um envelope de valores que deve ser
        medido pelos testes da cadeia. $\mathbb{Q}$ não é limitada a $16$
        bits; \code{Q16} é uma realização finita verificada de um envelope.
\end{enumerate}

%======================================================================
\section{Resultados já obtidos}\label{sec:ja}

\noindent\textbf{Objecto matemático versus realização finita.} O objecto
matemático é $\mathbb{Z}^{2}/{\sim}\cong\mathbb{Q}$, que é infinito. A
realização computacional \code{Qz}/\code{Q16} é uma janela finita medida
dentro desse objecto. Os testes verificam a janela, não a infinitude de
$\mathbb{Q}$.

\medskip\noindent
Cada realização computacional finita descrita neste paper já foi
\textbf{implementada} e \textbf{medida} nos domínios indicados. Os testes abaixo não substituem as demonstrações matemáticas:
verificam a realização computacional das identidades em domínios finitos
explicitamente especificados.
Os três ficheiros centrais são \code{lib/racionais.h},
\code{lib/dual16.h}, e \code{tests/racionais\_fixos.c}; a firma da torre
(\code{lib/torre\_alg.h}, \code{tests/torre\_alg.c}) amarra involução,
retração e subida ao mesmo degrau.

%----------------------------------------------------------------------
\subsection*{O par como tipo: \code{lib/racionais.h}}

As operações matemáticas preservam o domínio $\mathbb{Z}^{2}$. O envelope
canónico $E_{16}$ continua o \textbf{detector} do Gato:
\[
E_{16}=\bigl\{(p,q)\in\mathbb{Z}^{2}:
-32767\le p\le 32767,\;1\le q\le 32767,\;\gcd(|p|,q)=1\bigr\}.
\]
Hot path (20/08, \textsc{saturo}$\to$\textsc{promove}): fora de $E_{16}$,
\code{qz()} \emph{conta} (\code{qz\_saturou}) e \textbf{guarda} o representante
exacto em \code{int64} --- já não devolve $\pm32767$. \code{QzP} (I128) é o
testemunho acima do \code{int64}.
\begin{verbatim}
    typedef struct { int64_t p, q; } Qz;   /* E16 ou promovido; q != 0 */
    typedef struct { I128 p, q; } QzP;     /* fora do int64 */
\end{verbatim}
\code{Qz} não realiza todo $\mathbb{Q}$; realiza a classe exacta enquanto
couber em \code{int64}, com $E_{16}$ a marcar a primeira promoção.
$\operatorname{red}(0,b)=(0,1)$ para todo $b\neq 0$.

As operações da Def.~\ref{def:ops} estão realizadas:

\begin{center}\small
\begin{tabular}{@{}llll@{}}
\toprule
operação (§\ref{sec:ops}) & função & o que faz & nota \\
\midrule
soma        & \code{qz\_soma}     & $(ad'+cb',\,g\,b'd')$ com $g=\gcd(b,d)$, $b'=b/g$, $d'=d/g$ & cancela em cruz \\
produto     & \code{qz\_mult}     & $(ac',\,bd')$ com cancelamento & $\gcd(a,d)$ e $\gcd(c,b)$ \\
oposto      & \code{qz\_oposto}   & $(-a,\,b)$ & Lei 1 \\
inverso     & \code{qz\_inverso}  & $(b,\,a)$ & swap (involução multiplicativa) \\
divisão     & \code{qz\_divide}   & $a \otimes b^{-1}$ & recusa $b=0$ \\
redução     & \code{qz}           & divide por $\gcd$, $q>0$ (inverte ambos se necessário) & forma canónica do representante \\
\bottomrule
\end{tabular}
\end{center}

\noindent A saída de $E_{16}$ (\code{qz\_saturou}) é \textbf{contada à parte dos
defeitos}: o que não cabe no detector promove-se, e não se devolve tecto falso.

%----------------------------------------------------------------------
\subsection*{Produto exacto: par opera com par e resulta em par}

As operações deste degrau operam \textbf{pares com pares e resultam em pares}.
Cada produto inteiro entre componentes de $16$~bits é representado exactamente
por uma palavra de $32$~bits, materializada como dois blocos de $16$~bits
$(alto,baixo)$ --- sem depender de um tipo inteiro nativo de largura superior:
\begin{verbatim}
    typedef struct { uint16_t alto, baixo; } D32;
\end{verbatim}
A largura dupla é realizada estruturalmente por dois blocos de $16$~bits.
O par de palavras $(alto,baixo)$ representa exactamente os $32$~bits do
resultado em complemento de dois; a interpretação assinada ocorre após a
recomposição. Não se interpreta o overflow de um
\code{int16\_t} como o produto inteiro: a decomposição $(alto,baixo)$
\textbf{representa o resultado completo}. A redução ao domínio \code{Q16}
ocorre posteriormente, sob o envelope medido.

Todas as operações preservam a representação racional por pares; os produtos
intermediários podem exigir representação de largura dupla, aqui realizada por
duas palavras de $16$~bits:
\begin{itemize}
\item \textbf{soma}: $(a,b)\oplus(c,d)=(ad+bc,\;bd)$ --- os produtos
  intermediários são de largura dupla, o resultado é um par racional;
\item \textbf{produto}: $(a,b)\otimes(c,d)=(ac,\;bd)$ --- cada componente é
  um produto de $16$~bits em largura dupla, reduzido após cancelamento;
\item \textbf{comparação}: $\operatorname{sinal}(ad-bc)$ --- os produtos
  cruzados são de largura dupla, o sinal é lido do alto.
\end{itemize}

A comparação de ordem opera directamente nos pares:
\begin{verbatim}
    int d16_cmp_prod(int16_t a, int16_t b,
                     int16_t c, int16_t d);  /* sinal de a*b - c*d */
\end{verbatim}
que compara $ab$ com $cd$ \textbf{sem materializar nenhum dos dois como objecto
\code{Qz}}: compara o alto (interpretado como signed de $16$~bits) primeiro; em caso de
igualdade, compara o baixo (unsigned) --- produto cruzado exacto na
representação em palavras de $16$~bits.

A comparação $|a/b - c/d| < \varepsilon$, com $\varepsilon=e/f$, $f>0$
(\code{qz\_dist\_menor}) é reduzida a comparação inteira e opera inteiramente
por pares. Nenhuma das medições listadas utiliza \code{double}
nem decisões por limiar de ponto flutuante.

%----------------------------------------------------------------------
\subsection*{Dir e Cruz medidos: \code{tests/racionais\_fixos.c}}

O teste §Q4 varre $12^4 = 20\,736$ quádruplos ($a,b,c,d\in\{-6,\ldots,5\}$,
excluindo denominadores nulos) e verifica que a
\textbf{implementação computacional} da equivalência coincide entrada a entrada
com o predicado inteiro $ad=bc$ --- zero divergências.

\subsection*{Compatibilidade com o andar quadrático}

Os testes Q5--Q8 não pertencem ao núcleo deste degrau ($\mathbb{Z}\to\mathbb{Q}$);
verificam a interface com a realização já existente no andar superior.

\medskip\noindent
O teste §Q5 mede Dir e Cruz em cinco corpos quadráticos ($D=5,8,13,20,29$):
\begin{itemize}
  \item Dir tem parte $\sqrt D$ nula --- \textbf{é racional}.
  \item Cruz tem parte inteira nula --- \textbf{é irracional puro}.
  \item Dir $+$ Cruz devolve $x$ sem resto.
  \item Com a identidade no lugar da dobra, \emph{todos} seriam fixos --- o gume.
\end{itemize}
\textbf{Prova.} Para todo $D$ não quadrado, a conjugação de
$\mathbb{Q}(\sqrt{D})$ tem auto-espaço $+1$ igual a $\mathbb{Q}$.
\textbf{Verifica.} Q5 testa essa identidade nas instâncias
$D=5,8,13,20,29$.

O teste §Q7 mede a descida: $F(p,q)\neq 0$ em toda truncatura.

O teste §Q8 verifica o teorema dos resíduos: um resíduo sozinho é irracional,
e a soma dos dois cai em $\mathbb{Z}$ --- são os dois que dão o inteiro.

%----------------------------------------------------------------------
\subsection*{Firma da torre: \code{tests/torre\_alg.c}}

\noindent Este degrau $\mathbb{Z}\to\mathbb{Q}$ é o primeiro andar em que pares,
dualidade e produto cruzado coexistem na representação. O \emph{Corpo Algébrico}
(\code{algebrico thm:torre}) dobra a estrutura por
$T_{k+1}=T_k+T_k^{*}$; a \code{algebrico def:cone} separa
\textbf{involução} ($\nu\circ\nu=\mathrm{id}$) de \textbf{retração}
($\Sigma\circ\Pi=\mathrm{Id}$). O contrato computacional
\code{lib/torre\_alg.h} fixa três \emph{pipelines} sobre a mesma família de
operações --- involução, retração e subida --- sem confundi-los. A realização
passa por \code{lib/dual16.h}, \code{lib/dual32.h},
\code{lib/rt\_cf\_slot.h} e persistência \code{.mem}; o medidor §T1--§T4 verifica
que cada pipeline fecha no seu domínio e que
$\text{involução}\neq\text{retração}\neq\text{subida}$ mesmo quando, em casos
particulares, dois caminhos devolvem o mesmo valor numérico.
A síntese transversal destes três pipelines, slots e descida está em
\code{papers/arquitetura.tex}.

%----------------------------------------------------------------------
\subsection*{Inventário}

\begin{center}\small
\begin{tabular}{@{}llll@{}}
\toprule
afirmação & função / teste & ficheiro & resultado \\
\midrule
o par $(a,b)$ & \code{Qz \{p, q\}} & \code{lib/racionais.h} & tipo \\
soma por inteiros & \code{qz\_soma} & \code{lib/racionais.h} & $\checkmark$ \\
produto por inteiros & \code{qz\_mult} & \code{lib/racionais.h} & $\checkmark$ \\
inverso = swap & \code{qz\_inverso} & \code{lib/racionais.h} & $\checkmark$ \\
oposto = Lei 1 & \code{qz\_oposto} & \code{lib/racionais.h} & $\checkmark$ \\
redução pelo $\gcd$ & \code{qz} & \code{lib/racionais.h} & $\checkmark$ \\
saturação contada & \code{qz\_saturou} & \code{lib/racionais.h} & $\checkmark$ \\
$16\times16$ exacto em par & \code{d16\_mult} & \code{lib/dual16.h} & $\checkmark$ \\
ordem por prod.\ cruzado & \code{d16\_cmp\_prod} & \code{lib/dual16.h} & $\checkmark$ \\
produto $16\times16$ vs int32 & §E0--§EW & \code{tests/dual16.c} & $6/6$ \\
migração Q em $E_{16}$ & §GW, §G0--§G5 & \code{tests/migracao.c} & $7/7$ \\
descida $32\to16$ & §K2b & \code{tests/cadeia.c} & $\checkmark$ \\
RtOp coord.\ int16 & §R27 §T21 & \code{tests/reta.c} & $\checkmark$ \\
$|a-b|<\varepsilon$ & \code{qz\_dist\_menor} & \code{lib/racionais.h} & $\checkmark$ \\
$x^2$ vs $c$ & \code{qz\_cmp\_quad} & \code{lib/racionais.h} & $\checkmark$ \\
traço = inteiro & §Q1 & \code{tests/racionais\_fixos.c} & $60/60$ \\
norma $=-1$ (gume) & §Q2 & \code{tests/racionais\_fixos.c} & $60/60$ \\
$\sigma_m$ irracional & §Q3 & \code{tests/racionais\_fixos.c} & $60/60$ \\
equiv.\ por traço = $ad{=}bc$ & §Q4 & \code{tests/racionais\_fixos.c} & $20736/20736$ \\
$\operatorname{Fix}(\tau)=\mathbb{Q}$ em $\mathbb{Q}(\sqrt D)$ & §Q5 & \code{tests/racionais\_fixos.c} & $845$ \\
recorrência do pto fixo & §Q6 & \code{tests/racionais\_fixos.c} & $144/144$ \\
descida: $F\neq0$ & §Q7 & \code{tests/racionais\_fixos.c} & $\checkmark$ \\
resíduos $\to$ traço & §Q8 & \code{tests/racionais\_fixos.c} & $\checkmark$ \\
\midrule
\multicolumn{4}{@{}l@{}}{\emph{Firma da torre --- três pipelines (contrato \code{lib/torre\_alg.h}):}} \\
involução $\nu^2=\mathrm{id}$ & §T1 & \code{tests/torre\_alg.c} & $4/4$ \\
retração $\Sigma\circ\Pi=\mathrm{Id}$ & §T2 & \code{tests/torre\_alg.c} & $4/4$ \\
subida $T_k\to T_{k+1}$ & §T3 & \code{tests/torre\_alg.c} & $4/4$ \\
$\text{involução}\neq\text{retração}\neq\text{subida}$ & §T4 & \code{tests/torre\_alg.c} & $4/4$ \\
\bottomrule
\end{tabular}
\end{center}

\noindent \textbf{Nenhuma destas medições usa \code{double} ou limiar de ponto
flutuante}, nos caminhos de execução cobertos pelos testes Q1--Q8.

\medskip\noindent A construção de $\mathbb{Q}$ é o primeiro andar em que
pares, dualidade, produto cruzado e aritmética inteira aparecem
simultaneamente na representação do objecto; o andar seguinte muda de natureza porque introduz completude.
A migração para \code{int16} não está a inventar uma representação
computacional nova: está a materializar o andar
$\mathbb{Z}\to\operatorname{Frac}(\mathbb{Z})$ que já estava implícito na
torre.

%======================================================================
\section*{Síntese --- núcleo $\mathbb{Z}\!\to\!\mathbb{Q}$}

\[
\boxed{
\begin{array}{c}
\mathbb{Z}\xrightarrow{\text{duplicar}}V_{\mathrm{rat}}\subset\mathbb{Z}^{2}
\xrightarrow{\;q\;} V_{\mathrm{rat}}/{\sim}\;\cong\;\mathbb{Q}\\[2pt]
(a,b)\sim(c,d)\iff ad=bc,\qquad
\Psi([(a,b)])=a/b
\end{array}}
\]
\[
\boxed{
[(a,b)]\xrightarrow{\ \bar\iota\ }[(b,a)],\qquad
\Psi(\bar\iota(x))=1/\Psi(x)\ \ (x\neq 0),\qquad
[(a,b)]\otimes[(b,a)]\sim[(1,1)].
}
\]
Descida para 16 bits: equivalência por produto cruzado exacto em int32;
saturação contada à parte na escrita (\S\ref{sec:ja}).

\medskip\noindent
\code{python3 tools/refcruz.py} (chave \code{racionais}) ·
\code{inteiros.tex} · \code{corpo\_algebrico.tex} · \code{arquitetura.tex}.

\end{document}
